\documentclass[11pt]{article}
\usepackage{amsmath,amssymb,amsthm}
\usepackage[margin=1in]{geometry}
\usepackage{booktabs}
\usepackage[hidelinks]{hyperref}

\input{mef_macros.tex}

\numberwithin{equation}{section}

\theoremstyle{plain}
\newtheorem{theorem}{Theorem}[section]
\newtheorem{proposition}[theorem]{Proposition}
\newtheorem{lemma}[theorem]{Lemma}
\newtheorem{corollary}[theorem]{Corollary}
\theoremstyle{definition}
\newtheorem{definition}[theorem]{Definition}
\newtheorem{axiom}[theorem]{Axiom}
\newtheorem{problem}[theorem]{Problem}
\theoremstyle{remark}
\newtheorem{remark}[theorem]{Remark}

%% Local shorthands (no redefinition of mef_macros entries)
\newcommand{\M}{\mathcal{M}}
\newcommand{\Kr}{\mathcal{K}}
\newcommand{\Ric}{\operatorname{Ric}}
\newcommand{\II}{\mathrm{II}}
\newcommand{\N}{\mathcal{N}}
\newcommand{\A}{\mathcal{A}}
\newcommand{\Rr}{\mathcal{R}}
\newcommand{\Smatch}{\Sigma_{\mathrm{m}}}
\newcommand{\Geff}{G}

\title{The Schwarzschild interior from a smooth\\ twelve-dimensional
geometry}
\author{Dhiren J. Master\\[2pt]
\normalsize Independent Researcher\\[1pt]
\normalsize \texttt{dhiren@masterequation.org}}
\date{}

\begin{document}
\maketitle

\begin{abstract}
\noindent
The interior of a Schwarzschild black hole is a Kantowski--Sachs
cosmology whose singular asymptotics carry the Kasner data
$(-\tfrac13, \tfrac23, \tfrac23)$; a massless scalar field deforms
these data along a one-parameter family, the marginal direction of the
generalised Kasner constraint, and in four dimensions this amplitude is
free data. We study the system as the projection of a
twelve-dimensional geometry --- the four-dimensional chart lifted
against a fixed compact factor, carrying a smooth interface of finite
thickness --- subject to one structural axiom: the twelve-dimensional
metric is smooth. For matter satisfying the dominant energy condition
we prove the pointwise bound
$\Kr_{12} \geq (2592\pi^{2}/275)\, \Geff^{2}\rho^{2}$, so the axiom
excludes every bulk stress diverging towards the singular locus; matter
data survive only on the interface. The marginal scalar amplitude is
excluded with them: the constraint is exactly saturated, and the
interior data are determined by the mass alone. In every vacuum region
the four-dimensional metric is exactly Schwarzschild; the divergence of
its curvature is carried extrinsically, by a winding immersion in the
smooth ambient geometry. The global geometry joining interior to
exterior is posed as a two-point problem and reduced to five stated
estimates; the algebraic core of each step is machine-verified in
Lean~4. The Kasner triple is thereby recovered, not assumed.
\end{abstract}

\section{Introduction}\label{sec:intro}

The region inside the horizon of a Schwarzschild black hole is a
spatially homogeneous, anisotropic cosmology. In the interior the
Schwarzschild radial coordinate is timelike, the metric takes
Kantowski--Sachs form \cite{KantowskiSachs1966}, and the approach to
the singular locus $r \to 0$ is a vacuum Kasner regime: writing $\tau$
for proper time along the collapse, the metric approaches
\begin{equation}\label{eq:kasner4}
g_{4} \;\simeq\; -\,d\tau^{2} \;+\; \tau^{-2/3}\,dt^{2}
\;+\; \tau^{4/3}\, r_{0}^{2}\, d\Omega_{2}^{2}\,,
\end{equation}
the Kasner exponents $(-\tfrac13, \tfrac23, \tfrac23)$ being the unique
spherically symmetric vacuum data. When a massless scalar field is
present this rigidity is lost. The generalised Kasner constraint
acquires a source term, the exponents move along a one-parameter
family, and --- this is the point of departure for the present paper
--- the scalar amplitude is genuinely free data of the
four-dimensional problem: the Fuchsian analysis of quiescent
singularities \cite{AnderssonRendall2001, KichenassamyRendall1998}
constructs a solution for every choice of the amplitude, and nothing
in four-dimensional general relativity selects among them. The
behaviour is marginal in the precise sense that the scalar energy
density scales as $\tau^{-2}$, the same power as the anisotropy terms
in the Hamiltonian constraint, so the deformation neither decays nor
dominates but persists to the singularity. In higher-dimensional
vacuum gravity the same non-oscillatory, velocity-dominated regime is
available \cite{DemaretHenneauxSpindel1985, DamourHenneauxNicolai2003},
and the constraint structure is identical in form. To our knowledge,
no selection principle for the quiescent amplitude has appeared in the
literature.

The purpose of this paper is threefold.

(1) We consider the four-dimensional collapse as the projection of a
twelve-dimensional geometry: the chart is immersed in a product of a
Lorentzian block with a fixed compact factor, subject to a single
structural axiom --- the twelve-dimensional metric is smooth
(Axiom~\ref{ax:smooth}) --- and carrying a smooth interface of finite
thickness whose tension is a fixed input. For matter with
non-negative energy density and non-negative pressure sum (satisfied,
in particular, by the massless scalar and by every fluid considered
here, and implied for these matter forms by the dominant energy
condition) we prove a pointwise curvature bound: in dimension twelve,
\begin{equation}\label{eq:bound-intro}
\Kr_{12} \;\geq\; \frac{2592\,\pi^{2}}{275}\; \Geff^{2}\,
\rho_{\mathrm{tot}}^{2}
\end{equation}
(Lemma~\ref{lem:transmission}). The bound is algebraic and requires no
evolution; its consequence is Theorem~\ref{thm:exclusion}: under the
axiom, no bulk stress component whose energy density diverges towards
the singular locus is admissible. Matter data survive only on the
interface.

(2) The marginal scalar amplitude is excluded with the rest
(Corollary~\ref{cor:saturation}): the generalised Kasner constraint is
exactly saturated, the exponents are forced back to
$(-\tfrac13, \tfrac23, \tfrac23)$ with eight frozen internal
directions, and the interior data are determined by the mass alone.
In every vacuum region the four-dimensional chart metric is exactly
Schwarzschild (Proposition~\ref{prop:skeleton}); the divergence of the
four-dimensional curvature at the singular locus is carried
extrinsically, by a winding immersion of the chart in the smooth
ambient geometry, at the rate fixed by the Gauss equation
(Section~\ref{sec:skeleton}).

(3) The global geometry joining the winding interior description to
the static exterior product is posed as a two-point boundary problem
and reduced, by the method of matched asymptotic expansions, to five
explicitly stated estimates (Section~\ref{sec:twopoint}); existence of
the global interpolant is not claimed here. The algebraic core of each
completed step --- the saturation arithmetic and the coefficients of
the curvature bound --- is machine-verified in Lean~4
(Appendix~\ref{app:lean}).

The paper is organised as follows. Section~\ref{sec:setting} fixes the
twelve-dimensional setting, the axiom, the immersion of the chart, and
the reduced field equations, and derives the generalised Kasner
constraint in the form used throughout. Section~\ref{sec:bound} proves
the curvature bound. Section~\ref{sec:exclusion} proves the exclusion
theorem and its corollaries, including the saturation of the
constraint and the closure of the interface-absorption loophole.
Section~\ref{sec:skeleton} establishes the exact Schwarzschild
skeleton and describes the winding immersion. Section~\ref{sec:twopoint}
poses the two-point problem, states the five estimates, and records
the layered extension to the full inhomogeneous system.
Section~\ref{sec:discussion} discusses the scope of the results.
Appendix~\ref{app:lean} describes the machine verification.

\section{The setting}\label{sec:setting}

\subsection{The ambient geometry and the axiom}

Let $\mathrm{Sol}$ denote the three-dimensional mapping torus of a
hyperbolic monodromy: on coordinates $(x_{+}, x_{-}, z)$ the metric is
\begin{equation}\label{eq:sol}
g_{\mathrm{Sol}} \;=\; e^{2z}\,dx_{+}^{2} \;+\; e^{-2z}\,dx_{-}^{2}
\;+\; dz^{2}\,,
\end{equation}
compactified by $(x, z) \sim (Ax,\, z - L)$ for a fixed
$A \in SL(2,\mathbb{Z})$ hyperbolic with expanding eigenvalue
$e^{L} > 1$ along $x_{+}$; the scalar curvature is the constant $-2$.
Let $S^{3}$ carry the round metric of radius $R_{3}$, with polar angle
$\chi$ and latitude two-spheres of areal radius $R_{3}\sin\chi$, and
let $F$ be a fixed closed five-manifold of bounded geometry. The
ambient space is
\begin{equation}\label{eq:block}
\M_{12} \;=\; \mathbb{R}_{T} \times \mathrm{Sol} \times S^{3} \times
F\,, \qquad
g_{12} \;=\; -\,dT^{2} + g_{\mathrm{Sol}} + R_{3}^{2}\,d\Omega_{3}^{2}
+ g_{F}\,,
\end{equation}
a smooth homogeneous Lorentzian metric of bounded curvature. The field
equation is the twelve-dimensional Einstein equation
\begin{equation}\label{eq:einstein}
G_{MN} \;=\; 8\pi \Geff\, T_{MN} \;+\; \Lambda\, g_{MN}\,,
\end{equation}
with $T_{MN}$ comprising a canonical massless scalar field $\Phi$,
\begin{equation}\label{eq:scalar}
T^{\Phi}_{MN} \;=\; \nabla_{M}\Phi\,\nabla_{N}\Phi
- \tfrac12\, g_{MN}\,\lvert\nabla\Phi\rvert^{2}\,, \qquad
\Box\,\Phi = 0\,,
\end{equation}
together with a perfect-fluid sector, and $\Lambda$ a cosmological
constant whose contributions are booked and uniformly subleading at
the curvature scales considered; we set $\Lambda = 0$ in the estimates
and record its effect in Remark~\ref{rem:lambda}. A domain wall ---
the interface --- separates two field sectors of the theory; it is a
smooth localised profile of finite thickness $\epsilon$ in one
internal direction, of fixed tension $\sigma$, and the thin-interface
(Israel \cite{Israel1966}) treatment used below is the computational
idealisation $\epsilon \to 0$ of this smooth object.

\begin{axiom}[Smoothness]\label{ax:smooth}
Admissible configurations are smooth solutions of
\eqref{eq:einstein} on $\M_{12}$; equivalently, the twelve-dimensional
Kretschmann scalar $\Kr_{12} = R_{MNPQ}R^{MNPQ}$ is locally bounded.
The interface is the smooth profile above; statements referring to
quantities ``off the interface'' are bookkeeping of the thin-interface
limit, not an exemption from the axiom.
\end{axiom}

One should be careful about the interpretation of
Axiom~\ref{ax:smooth}. It is a definition of the solution class, not a
theorem about it: configurations violating it are excluded as
unphysical, in the same way that geodesically incomplete extensions
are excluded in the classical setting. Every result below is
conditional on this class, and the question of whether the class is
non-empty for the global problem --- the existence of the interpolant
of Section~\ref{sec:twopoint} --- is open and is stated as such. The
axiom, a property of its members, and the existence of a member are
three separate matters, and the paper establishes only the second.

\subsection{The immersed chart and the reduced system}

The four-dimensional chart is the immersion
\begin{equation}\label{eq:map}
X : (\tau, r, \theta, \phi) \;\longmapsto\;
\bigl(\, T(\tau, r),\;\; x_{+} = r/\ell,\;\; z(\tau, r),\;\;
\chi(\tau, r),\;\; \theta,\; \phi \,\bigr) \;\subset\; \M_{12}\,,
\end{equation}
at fixed $x_{-}$ and fixed position on $F$, with $\ell = e^{z_{0}}$ a
fixed normalisation; the chart is spherically symmetric, with
two-dimensional orbit space parametrised by $(\tau, r)$, and has
codimension eight. Pulling back \eqref{eq:block} along \eqref{eq:map}
gives the induced metric
\begin{equation}\label{eq:induced}
g_{4} \;=\; -\,\N^{2}\,d\tau^{2} + 2\,\mathcal{S}\,d\tau\,dr
+ \A^{2}\,dr^{2} + \Rr^{2}\,d\Omega_{2}^{2}
\end{equation}
with the exact dictionary
\begin{equation}\label{eq:dictionary}
\begin{aligned}
\N^{2} &= \dot T^{2} - \dot z^{2} - R_{3}^{2}\dot\chi^{2}\,, &
\mathcal{S} &= -\dot T T' + \dot z z' + R_{3}^{2}\dot\chi\chi'\,,\\
\A^{2} &= e^{2(z - z_{0})} - T'^{2} + z'^{2} + R_{3}^{2}\chi'^{2}\,, &
\R &= R_{3}\sin\chi\,,
\end{aligned}
\end{equation}
where dots and primes denote $\partial_{\tau}$ and $\partial_{r}$.
The reduced field equations on the chart are the projections of
\eqref{eq:einstein}: the doubly contracted Gauss equation
\begin{equation}\label{eq:gauss}
G_{ab}[g_{4}] \;=\; 8\pi\Geff\, T_{ab}
+ \mathcal{E}_{ab}[\II] + \bar{\mathcal{C}}_{ab}\,,
\end{equation}
in which $\mathcal{E}_{ab}[\II]$ is the standard quadratic in the
second fundamental form over the eight normals and
$\bar{\mathcal{C}}_{ab}$ collects the tangential projections of the
ambient curvature, bounded by homogeneity of \eqref{eq:block}; the
Codazzi equations in the two non-trivial normals; the thin-interface
junction condition at tension $\sigma$; and the scalar equation
$\Box\Phi = 0$. The full component expansion of \eqref{eq:gauss} is
lengthy and is generated symbolically in the machine-verification
pipeline of Appendix~\ref{app:lean} rather than transcribed; what the
arguments below use is the dictionary \eqref{eq:dictionary}, the
constraint structure, and the near-singular limit derived next.

\subsection{The generalised Kasner constraint}\label{sec:budget}

In the near-singular regime the chart is homogeneous
($T' = z' = \chi' = 0$), velocity-dominated, and the induced metric
takes the diagonal form
$-d\tau^{2} + \sum_{i=1}^{11}\tau^{2p_{i}}(\sigma^{i})^{2}$ on the
eleven spatial directions: three along the chart (the $x_{+}$ fibre
and the latitude two-sphere) and eight frozen normals. With
$\Phi = \Phi_{0} + q\ln\tau$ the $(\tau\tau)$-component and trace of
\eqref{eq:einstein} give, at leading order,
\begin{equation}\label{eq:budget}
\sum_{i=1}^{11} p_{i} \;=\; 1\,, \qquad
\sum_{i=1}^{11} p_{i}^{2} \;=\; 1 \;-\; 8\pi\Geff\, q^{2}\,,
\end{equation}
the higher-dimensional generalised Kasner constraint
\cite{DemaretHenneauxSpindel1985}; a stiff fluid component
$\rho = s\,\tau^{-2}$ enters identically with $8\pi\Geff q^{2}$
replaced by $16\pi\Geff s$. Each amplitude tilts the exponents along
the solution family of \eqref{eq:budget}, and every member of the
family is realised by a quiescent solution in four dimensions
\cite{AnderssonRendall2001}. The point is here that the deformation is
marginal: the source scales as $\tau^{-2}$, exactly the power of the
anisotropy terms, so the amplitude neither decays nor dominates ---
it tilts the exponents permanently, by an amount the constraint
records, and in the four-dimensional problem nothing fixes it
\cite{AnderssonRendall2001}. The vacuum data
\begin{equation}\label{eq:winding-data}
(p_{1}, \dots, p_{11}) \;=\;
\bigl(-\tfrac13,\ \tfrac23,\ \tfrac23,\ 0, \dots, 0\bigr)
\end{equation}
saturate both sums exactly; the saturation arithmetic, the equivalence
of the constraint forms, and the statement that saturation forces
every internal exponent and the source amplitude to vanish termwise
are machine-verified (Appendix~\ref{app:lean}).

\section{The curvature bound}\label{sec:bound}

\begin{lemma}[Pointwise transmission]\label{lem:transmission}
Let $g$ satisfy \eqref{eq:einstein} with $\Lambda = 0$ at a point of a
$D$-dimensional Lorentzian manifold, $D \geq 4$, and suppose the
source has energy density $\rho_{\mathrm{tot}} \geq 0$ and pressure
sum $\sum_{i} p_{i} \geq 0$ in an orthonormal frame with timelike leg
$u$. Then, at that point,
\begin{equation}\label{eq:bound-genD}
\Kr_{D} \;\geq\; \frac{2}{D-1}\,\bigl(R_{uu}\bigr)^{2}
\;\geq\; \frac{2}{D-1}\Bigl(\frac{8\pi\,(D-3)}{D-2}\Bigr)^{\!2}
\Geff^{2}\,\rho_{\mathrm{tot}}^{2}\,.
\end{equation}
In dimension twelve the constant evaluates to
\begin{equation}\label{eq:bound}
\Kr_{12} \;\geq\; \frac{2}{11}\,\bigl(R_{uu}\bigr)^{2}
\;\geq\; \frac{2592\,\pi^{2}}{275}\;\Geff^{2}\,
\rho_{\mathrm{tot}}^{2}\,.
\end{equation}
\end{lemma}

\begin{proof}
The Riemann tensor decomposes as Weyl plus Ricci parts, giving in
dimension $D$
\[
\Kr \;=\; \lvert C\rvert^{2} \;+\; \frac{4}{D-2}\,
\lvert\Ric\rvert^{2} \;-\; \frac{2}{(D-1)(D-2)}\,R^{2}\,,
\]
and $R^{2} \leq D\,\lvert\Ric\rvert^{2}$ by the Cauchy--Schwarz
inequality applied to $R = g^{MN}R_{MN}$. Since
$\lvert C\rvert^{2} \geq 0$,
\[
\Kr \;\geq\; \Bigl[\frac{4}{D-2} - \frac{2D}{(D-1)(D-2)}\Bigr]
\lvert\Ric\rvert^{2} \;=\; \frac{2}{D-1}\,\lvert\Ric\rvert^{2}
\;\geq\; \frac{2}{D-1}\,(R_{uu})^{2}\,,
\]
which is the first inequality of \eqref{eq:bound-genD}. Trace
reversal of \eqref{eq:einstein} gives
$R_{MN} = 8\pi\Geff\bigl(T_{MN} - \tfrac{1}{D-2}g_{MN}T\bigr)$, whence
\[
R_{uu} \;=\; 8\pi\Geff\Bigl[\rho_{\mathrm{tot}}\,\frac{D-3}{D-2}
\;+\; \frac{1}{D-2}\sum_{i}p_{i}\Bigr]
\;\geq\; \frac{8\pi\,(D-3)}{D-2}\,\Geff\,\rho_{\mathrm{tot}}\,,
\]
using both hypotheses; together these give \eqref{eq:bound-genD}. No
step is particular to any dimension. At $D = 12$ the two coefficients
are $\tfrac{2}{11}$ and $\tfrac{9}{10}$, and
$\tfrac{2}{11}\bigl(8\pi\cdot\tfrac{9}{10}\bigr)^{2} =
\tfrac{2592\pi^{2}}{275}$, which is \eqref{eq:bound}. The canonical scalar \eqref{eq:scalar}
satisfies both hypotheses, with $\rho_{\Phi} = p_{\Phi} \geq 0$ along
its gradient regime, and is covered by the same inequality. The
rational coefficients are machine-verified
(Appendix~\ref{app:lean}).
\end{proof}

The estimate is pointwise, algebraic, and frame-covariant; no
evolution enters, and no gauge choice affects it. Elementary and valid
in every dimension, it is perhaps of interest in its own right as a
constraint on bulk matter in any lifted collapse; its sole role in
this paper is Theorem~\ref{thm:exclusion}.

\begin{remark}\label{rem:lambda}
With $\Lambda \neq 0$ the trace reversal contributes
$2\Lambda/(D-2)$ to $R_{uu}$ and \eqref{eq:bound} acquires a cross
term of order $\Lambda\rho_{\mathrm{tot}}$; for the booked value of
$\Lambda$ this is uniformly negligible at every curvature scale
considered in this paper, and we do not display it further.
\end{remark}

\section{The exclusion theorem}\label{sec:exclusion}

\begin{theorem}[Bulk exclusion at the singular locus]
\label{thm:exclusion}
Under Axiom~\ref{ax:smooth}, no admissible solution carries a bulk
stress component whose energy density diverges towards the singular
locus. Consequently the ambient geometry is vacuum in a neighbourhood
of the singular locus off the interface: all matter data reside on the
interface, and the divergence of the four-dimensional chart curvature
is carried extrinsically, through $\mathcal{E}_{ab}[\II]$ in
\eqref{eq:gauss}.
\end{theorem}

\begin{proof}
Suppose a bulk component has
$\rho_{\mathrm{tot}}(\tau) \to \infty$ as $\tau \to 0$ on a
neighbourhood of ambient points off the interface. By
Lemma~\ref{lem:transmission}, $\Kr_{12} \geq
\tfrac{2592\pi^{2}}{275}\Geff^{2}\rho_{\mathrm{tot}}^{2} \to \infty$
there. Since $\tau$ is proper time along the collapse, the curvature
diverges along timelike curves at finite proper time; this is a
curvature singularity, hence a violation of Axiom~\ref{ax:smooth}
under either its curvature or its completeness reading. The estimate
being pointwise and frame-covariant, no coordinate choice evades it.
With every divergent bulk component excluded, the ambient stress is
bounded near the singular locus, the block \eqref{eq:block} retains
its bounded homogeneous curvature, and the only remaining carriers of
divergence are the extrinsic sector of the immersed chart and the
interface data.
\end{proof}

\begin{corollary}[Saturation]\label{cor:saturation}
The scalar and stiff amplitudes vanish at the singular locus:
$\rho_{\Phi} = \tfrac12 q^{2}\tau^{-2}$ and
$\rho_{\mathrm{stiff}} = s\,\tau^{-2}$ diverge for any $q \neq 0$,
$s > 0$ and are excluded by Theorem~\ref{thm:exclusion}. The
constraint \eqref{eq:budget} is therefore exactly saturated, the
exponents are \eqref{eq:winding-data}, and the interior data carry no
free function: they are determined by the mass alone.
\end{corollary}

\begin{proof}
Immediate from Theorem~\ref{thm:exclusion} and the termwise saturation
statement of Section~\ref{sec:budget}.
\end{proof}

\begin{corollary}\label{cor:allspecies}
The same exclusion applies to every species: a radiation component
($\rho \propto \tau^{-4/3}$), or any component with
$\rho \to \infty$, is excluded identically. The near-singular bulk is
exactly vacuum.
\end{corollary}

\begin{lemma}[No interface absorption]\label{lem:wall}
The excluded amplitudes cannot be relocated to the interface: a
localised stiff component of surface amplitude $s_{\mathrm{w}} > 0$
shifts the junction data away from the fixed tension,
$S_{ab} \to \sigma\, u_{a}u_{b} + s_{\mathrm{w}}(\cdots)$, and since
$\sigma$ is a single fixed input of the model with no residual surface
freedom, the exterior matching of Section~\ref{sec:twopoint} fails at
first order in $s_{\mathrm{w}}$. Hence $s_{\mathrm{w}} = 0$: the
interface carries exactly its fixed data.
\end{lemma}

We note that Corollary~\ref{cor:saturation} sits in a specific
relation to the four-dimensional theory. There, the scalar amplitude
is free data --- the quiescent solutions of
\cite{AnderssonRendall2001} exist for every $q$ --- and no
four-dimensional principle selects among them. On the lift the
selection is made by Axiom~\ref{ax:smooth} alone: the marginal
direction of the constraint is precisely the direction along which
bulk stress survives to the singular locus, and smoothness closes it.

\section{The exact skeleton and the winding immersion}
\label{sec:skeleton}

\begin{proposition}[Exact Schwarzschild skeleton]\label{prop:skeleton}
In every region where the bulk is vacuum, the four-dimensional chart
metric is exactly the Schwarzschild metric of a single mass $M$. In
particular, by Corollary~\ref{cor:allspecies} this holds in a full
neighbourhood of the singular locus, where the chart metric is the
Schwarzschild interior in Kantowski--Sachs form with the exponents
\eqref{eq:winding-data}; and it holds in the static exterior region.
Deviations from the single Schwarzschild skeleton are confined to the
matter-bearing intermediate region.
\end{proposition}

\begin{proof}
Spherically symmetric vacuum solutions of the four-dimensional
projected system with bounded extrinsic corrections reduce, in the
regions stated, to spherically symmetric vacuum solutions of the
Einstein equations, and these are locally Schwarzschild
\cite{Birkhoff1923}; near the singular locus the extrinsic and ambient
terms in \eqref{eq:gauss} enter the constraint at the orders booked in
Section~\ref{sec:budget}, and the vacuum statement is
Corollary~\ref{cor:allspecies}. Continuity of the Misner--Sharp mass
\cite{MisnerSharp1964} across the intermediate region identifies the
two vacuum masses.
\end{proof}

The four-dimensional curvature nevertheless diverges at the singular
locus, $\Kr_{4} = 48\,\Geff^{2}M^{2}/r^{6}$, while
Theorem~\ref{thm:exclusion} keeps $\Kr_{12}$ bounded there. The two
statements are reconciled by the immersion: on the data
\eqref{eq:winding-data} the chart is realised as
\begin{equation}\label{eq:winding}
z(\tau) \;=\; z_{0} - \tfrac13\ln\tau \ (\mathrm{mod}\ L)\,, \qquad
\sin\chi(\tau) \;=\; \frac{r_{0}}{R_{3}}\,\tau^{2/3}\,, \qquad
\dot T^{2} - \dot z^{2} - R_{3}^{2}\dot\chi^{2} \;=\; 1\,,
\end{equation}
which reproduces \eqref{eq:kasner4} exactly through the dictionary
\eqref{eq:dictionary}: the anisotropy factor $\tau^{-2/3}$ is realised
as motion along the compact Sol base, the immersion winding once
around the mapping torus per decrease of $z$ by $L$, so that the
winding number to time $\tau$ is $n(\tau) = \ln(1/\tau)/(3L)$, growing
without bound while the ambient geometry through which the chart winds
remains homogeneous and smooth. The divergence of $\Kr_{4}$ is carried
entirely by the second fundamental form, whose growth rate the Gauss
equation fixes against $\Kr_{4}^{1/2}$; nothing diverges in the
ambient geometry. The construction is a compact-direction analogue of
the classical isometric embeddings of the Schwarzschild geometry
\cite{Fronsdal1959}, with the unbounded stretch supplied by winding
rather than by non-compact hyperbolic coordinates. The classical
embeddings, however, employ a flat ambient space carrying no field
content; here the ambient geometry is itself a solution of
\eqref{eq:einstein}, and, to our knowledge, an embedding of the
Schwarzschild interior in which the curvature divergence is carried by
winding a compact direction of a dynamical background has not
previously appeared.

\section{The two-point problem}\label{sec:twopoint}

The exterior of the configuration is the static product
$\mathrm{Schwarzschild}(M) \times K$, with $K$ the fixed compact
factor: for a metric product the Ricci tensor is block diagonal, the
four-dimensional block is Ricci-flat, and the product solves
\eqref{eq:einstein} wherever $\Kr_{4}$ is bounded --- in particular
through a full neighbourhood of the horizon --- failing only at the
singular locus, which is precisely where the winding interior
description takes over. The two descriptions overlap at intermediate
curvature. An interpolating solution, if it exists, is a solution of
the reduced system on the orbit space joining the interior data of
Corollary~\ref{cor:saturation} --- which carry no free function ---
to the exterior product data across the overlap, with the interface at
its fixed tension.

\begin{problem}\label{prob:twopoint}
For given $M > 0$, find a solution of the reduced system
\eqref{eq:gauss} on the orbit space, in the diagonal gauge
$\mathcal{S} = 0$, with: the interior asymptotics
\eqref{eq:winding-data}--\eqref{eq:winding}; the exterior data
$(g_{4}, \II) = (\mathrm{Schwarzschild}(M), 0)$ plus the junction
condition at tension $\sigma$ on the matching surface; profiles $C^{1}$
across the horizon locus in a regular chart; and the winding sector
reconciled in the sense of (M3) below.
\end{problem}

Three matching layers arise, in ascending non-triviality. (M1)
\emph{Four-dimensional data:} by Proposition~\ref{prop:skeleton} the
vacuum sectors on both sides are the same Schwarzschild($M$), so the
leading-order four-dimensional matching is satisfied identically;
numerically, in a dimensionless implementation of the reduced system,
the mass functional at horizon crossing varies by less than one per
cent across two decades of the intermediate-region matter parameters.
(M2) \emph{Interface data:} the one-sided junction condition at
tension $\sigma$, a finite condition, not a free function. (M3)
\emph{The winding reconciliation:} the interior description winds,
$n(\tau) \to \infty$ towards the singular locus and quantised per Sol
period, while the exterior product is winding-static; an interpolant
must supply an admissible transition geometry on the overlap unwinding
the one description into the other. This layer is invisible to the
four-dimensional problem and is where the twelve-dimensional content
of Problem~\ref{prob:twopoint} resides.

By the method of matched asymptotic expansions the problem reduces to
a fixed-point formulation: the interior admits an expansion in powers
of $\tau^{2/3}$ about the exact asymptote (both correction channels,
the sphere curvature and any bounded matter entering from the
intermediate region, arise at this relative order), with no free
coefficient at the singular end, anchored by the Fuchsian theory of
velocity-dominated singularities \cite{KichenassamyRendall1998,
AnderssonRendall2001}; the exterior is exact; and a solution of
Problem~\ref{prob:twopoint} corresponds to a fixed point of the
remainder map on a weighted space over the overlap. We reduce the
existence question to five estimates, stated here and left open:
(EC1) remainder bounds for the interior expansion, uniform in $M$ on
compact ranges; (EC2) propagation of $C^{0}$ bounds through the
matter-bearing intermediate region; (EC3) uniform estimates across the
transversally crossed horizon locus in a regular chart; (EC4)
smallness of the Lipschitz constant of the remainder map on the
overlap; (EC5) construction of the transition geometry of (M3),
compatible with (M2). We pose the problem and reduce it; we do not
claim its solution. The transition geometry of (M3) has no
four-dimensional counterpart, and it would be interesting to construct
it explicitly, beginning with the quantised unwinding across a single
Sol period. The extension of Problem~\ref{prob:twopoint} from
the orbit-space reduction to the full inhomogeneous system is
structurally identical, with one additional requirement --- control of
the inhomogeneous perturbations by the reduction --- and is recorded
as the layered form of the problem.

\section{Discussion}\label{sec:discussion}

One should be careful about the scope of these results. Everything
proved here is conditional on the solution class of
Axiom~\ref{ax:smooth} and on the specific lift
\eqref{eq:block}--\eqref{eq:map}: the theorems say nothing about other
lifts, other dimensions, or the four-dimensional theory taken alone,
in which the scalar amplitude remains free data. Nothing in the
exclusion mechanism is specific to the dimension --- the bound
\eqref{eq:bound-genD} holds for every $D \geq 4$ --- and twelve is the
case under study. Existence of the
global interpolant is open; the paper reduces it to the five estimates
of Section~\ref{sec:twopoint} and stops. No modification of the
exterior Schwarzschild geometry is asserted at the level of the
reduction: Proposition~\ref{prop:skeleton} places the four-dimensional
metric exactly on the classical solution wherever the bulk is vacuum,
and the content of the construction is interior and
twelve-dimensional. Within these limits, the results admit a plain
summary. In four dimensions the singular asymptotics of spherical
collapse with a massless scalar form a one-parameter family, and the
selection among its members is undetermined. On the smooth
twelve-dimensional lift the selection is a theorem: the bound
\eqref{eq:bound} transmits any surviving bulk amplitude into a
curvature divergence, the axiom excludes it, and the vacuum Kasner
data \eqref{eq:winding-data} --- with the four-dimensional geometry
exactly Schwarzschild and the divergence carried by winding in a
smooth ambient space --- are what remains. The Kasner triple of the
Schwarzschild interior is thereby recovered, not assumed.

\section*{Acknowledgements}

The author used Claude (Anthropic) and Gemini (Google DeepMind) as
mathematical development and analytical tools during the preparation
of this work. All intellectual content, theoretical framework, and
scientific claims are the sole responsibility of the author.

\appendix

\section{Machine verification}\label{app:lean}

Two Lean~4 certificates (Mathlib; rational arithmetic, decidable
throughout) accompany the paper and have been compiled with all goals
passed. The first certifies the algebraic layer of
Section~\ref{sec:budget}: the saturation of \eqref{eq:budget} by the
data \eqref{eq:winding-data}; the factorisation of the
$\tau$-dependence out of the constraint; the equivalence of the raw
and canonical constraint forms; and the termwise statement that
saturation forces the eight internal exponents and the source
amplitude to vanish, the amplitude entering as a single non-negative
rational so that the certificate is independent of normalisation
conventions. The second certifies the rational layer of
Lemma~\ref{lem:transmission}: the coefficient identities
$4/(D-2) - 2D/((D-1)(D-2)) = 2/(D-1) = 2/11$ and
$(D-3)/(D-2) = 9/10$ at $D = 12$; the composite constant
$\tfrac{2}{11}\bigl(8\cdot\tfrac{9}{10}\bigr)^{2} =
\tfrac{2592}{275}$; monotonicity of the bound in the density; and the
exclusion dichotomy --- boundedness of $s/t^{2}$ for arbitrarily
small $t$ forces $s = 0$ --- by explicit rational witness, which is
the algebraic core of Corollary~\ref{cor:saturation}. In addition,
the component equations of the reduced system \eqref{eq:gauss} used in
the numerical remark of Section~\ref{sec:twopoint} are generated
symbolically from \eqref{eq:block} and \eqref{eq:map} by computer
algebra, with the induced-metric dictionary \eqref{eq:dictionary},
the vanishing of the second fundamental form in the spectator
normals, and the contracted Gauss identity verified within the same
pipeline, so that no hand-transcribed tensor computation enters the
numerics.

\begin{thebibliography}{99}

\bibitem{KantowskiSachs1966} R.~Kantowski and R.~K.~Sachs, Some
spatially homogeneous anisotropic relativistic cosmological models,
\emph{J. Math. Phys.} \textbf{7} (1966) 443--446.

\bibitem{BKL1970} V.~A.~Belinskii, I.~M.~Khalatnikov and
E.~M.~Lifshitz, Oscillatory approach to a singular point in the
relativistic cosmology, \emph{Adv. Phys.} \textbf{19} (1970)
525--573.

\bibitem{AnderssonRendall2001} L.~Andersson and A.~D.~Rendall,
Quiescent cosmological singularities, \emph{Commun. Math. Phys.}
\textbf{218} (2001) 479--511.

\bibitem{KichenassamyRendall1998} S.~Kichenassamy and A.~D.~Rendall,
Analytic description of singularities in Gowdy spacetimes,
\emph{Class. Quantum Grav.} \textbf{15} (1998) 1339--1355.

\bibitem{DemaretHenneauxSpindel1985} J.~Demaret, M.~Henneaux and
P.~Spindel, Non-oscillatory behaviour in vacuum Kaluza--Klein
cosmologies, \emph{Phys. Lett. B} \textbf{164} (1985) 27--30.

\bibitem{DamourHenneauxNicolai2003} T.~Damour, M.~Henneaux and
H.~Nicolai, Cosmological billiards, \emph{Class. Quantum Grav.}
\textbf{20} (2003) R145--R200.

\bibitem{Israel1966} W.~Israel, Singular hypersurfaces and thin
shells in general relativity, \emph{Nuovo Cimento B} \textbf{44}
(1966) 1--14.

\bibitem{MisnerSharp1964} C.~W.~Misner and D.~H.~Sharp, Relativistic
equations for adiabatic, spherically symmetric gravitational
collapse, \emph{Phys. Rev.} \textbf{136} (1964) B571--B576.

\bibitem{Birkhoff1923} G.~D.~Birkhoff, \emph{Relativity and Modern
Physics}, Harvard University Press, Cambridge MA, 1923.

\bibitem{Fronsdal1959} C.~Fronsdal, Completion and embedding of the
Schwarzschild solution, \emph{Phys. Rev.} \textbf{116} (1959)
778--781.

\bibitem{Wald1984} R.~M.~Wald, \emph{General Relativity}, University
of Chicago Press, Chicago, 1984.

\end{thebibliography}

\end{document}
